% ===== overview =====
\begin{frame}{이번 주차 개요}
  \begin{columns}[T]
    \begin{column}{0.58\textwidth}
      카지노의 슬롯머신: 각 손잡이(arm)는 \textbf{모르는} 평균 보상을
      가진 보상 기계. 제한된 횟수만 당길 수 있다면, 어떤 순서로
      당겨야 총 보상을 최대화할 수 있을까? --- 이 단순한 질문이
      \kb{멀티암 밴딧} 문제.
      \vskip0.4em
      \begin{itemize}
        \item 상태(state)도, 전이(transition)도 \textbf{없다} ---
          \kb{탐험과 활용의 딜레마} 하나만 순수하게 남긴
          \textbf{가장 작은 뼈대}.
        \item 다음 장(MDP)에서는 바로 이 밴딧에 ``상태''라는 재료
          하나를 더해서 수학적으로 정식화한다.
      \end{itemize}
      \vskip0.4em
      \begin{block}{이 주차를 마치면}
        $k$개 팔 밴딧을 \textbf{총보상·후회}로 정식화하고,
        ε-greedy·낙관적 초기화·UCB를 ``탐험이 어떻게 만들어지는가''
        기준으로 비교하며, 밴딧을 ``상태가 하나뿐인 MDP''로
        해석할 수 있다.
      \end{block}
    \end{column}
    \begin{column}{0.4\textwidth}
      \centering
      % [fig missing] \includegraphics[width=0.95\linewidth]{ch02_1_bandit_octopus.png}
      \vskip0.3em
      {\footnotesize 이름의 유래: \textit{bandit machine}(슬롯머신)
      + 팔이 여러 개인 \textbf{문어}의 말장난.}
      \vskip0.3em
      {\footnotesize 1930년대 말부터 수학 연구, 1952년
      Robbins의 후회 하한($O(\ln T)$)으로 확립.}
      \vskip0.3em
      {\footnotesize 두 잣대: \textbf{총보상} $\sum R_t$
      (최대화)·\textbf{후회}(오라클과의 차이, 최소화).}
    \end{column}
  \end{columns}
\end{frame}
